\uuid{PBKf}
\titre{Noyau/Image/Injection/Surjection/Bijection/Théorème du rang}
\auteur{hanani ; mzari}
\organisation{exo7}
\licence{CC BY-NC-SA 4.0}
\datecreate{2020-01-20}
\niveau{L1}
\module{Algèbre}
\chapitre{Application linéaire}
\sousChapitre{Définition}
\difficulte{2}
\theme{Noyau et image, Noyau, Image, Injection, Surjection, Bijection, Théorème du rang}
\exo7qcmid{lille-2.applications-lineaires.016}

\contenu{
\begin{question}
\qclassification{107.01}
\qtags{motcle=Noyau/Image/Injection/Surjection/Bijection/Théorème du rang}
On considère l' application linéaire :
$$\begin{array}{rccc}f:&\Rr^3&\to&\Rr^3\\
& (x,y,z)&\to &(x-y,y-z,x+z). \,  \end{array}$$
Quelles sont les assertions vraies ?
\begin{answers}
\bad{$\dim \ker f = 1$}.
\good{$f$ est injective}.
\good{$\dim \Im f = 3$}.
\bad{$f$ n'est pas bijective}.
\end{answers}
\begin{explanations} $\ker f = \{(x,y,z) \in \Rr^3 \; ; \; (x-y,y-z,x+z)=(0,0,0)\} = \{(0,0,0)\}$. Donc $\dim \ker f = 0$ et  $f$ est injective.
\vskip0mm
D'après le théorème du rang, $\dim \Im f = 3=\dim \Rr^3$ et comme $\Im f$ est un sous-espace de $\Rr^3$, $\Im f=\Rr^3$, donc $f$ est surjective. Par conséquent, $f$ est bijective.
\end{explanations}
\end{question}
}
